\documentclass[10pt]{article}
\usepackage[T1]{fontenc}
\usepackage[utf8]{inputenc}
\usepackage[french]{babel}
\usepackage{amsmath,amssymb,mathrsfs,mathpazo}
\usepackage{geometry}
\geometry{a4paper,margin=2cm}
\usepackage{graphicx,array,tabularx,tabularray,colortbl,multirow,multicol,booktabs,enumitem,pifont,ulem,fancybox,diagbox,cancel}
\usepackage{pstricks,pstricks-add,pst-plot,pst-tree,pst-node,pst-text,pst-func,pst-math,pst-bspline,pst-3dplot,pst-eucl,pst-solides3d}
\let\pscalculate\undefined
\usepackage{pst-bezier}
\usepackage{tikz,pgfplots,tkz-tab}
\usetikzlibrary{arrows,arrows.meta,patterns,calc,positioning,trees,decorations.markings,shapes.geometric,shapes.misc,shapes.arrows,intersections,decorations.pathmorphing,decorations.pathreplacing,backgrounds,shadows,babel}
\pgfplotsset{compat=1.18}
\usepgfplotslibrary{fillbetween}
\usepackage[np]{numprint}
\usepackage{listings,makecell,pas-tableur,twoopt}
\let\barre\undefined
\usepackage[pstricks]{tablvar}
\definecolor{darkgreen}{RGB}{0,100,0}
\providecommand{\red}{\color{red}}
\providecommand{\blue}{\color{blue}}
\providecommand{\green}{\color{green}}
\providecommand{\black}{\color{black}}
\providecommand{\white}{\color{white}}
\providecommand{\gray}{\color{gray}}
\providecommand{\cyan}{\color{cyan}}
\providecommand{\magenta}{\color{magenta}}
\providecommand{\yellow}{\color{yellow}}
\providecommand{\violet}{\color{violet}}
\providecommand{\orange}{\color{orange}}
\providecommand{\euro}{\texteuro}
\providecommand{\EUR}{\texteuro}
\providecommand{\ieme}{\textsuperscript{e}}
\providecommand{\degres}{\ensuremath{^\circ}}
\providecommand{\degree}{\ensuremath{^\circ}}
\providecommand{\vv}[1]{\overrightarrow{#1}}
\providecommand{\Coord}[1]{#1}
\providecommand{\decosix}{}
\providecommand{\starredbullet}{\ensuremath{\bullet}}
\pagestyle{empty}
\setlength{\parindent}{0pt}
\usepackage{adjustbox}\DeclareUnicodeCharacter{2236}{:}\DeclareUnicodeCharacter{2190}{\ensuremath{\leftarrow}}
\providecommand{\R}{\mathbb{R}}
\providecommand{\N}{\mathbb{N}}
\providecommand{\Z}{\mathbb{Z}}
\providecommand{\Q}{\mathbb{Q}}
\providecommand{\D}{\mathbb{D}}
\providecommand{\C}{\mathbb{C}}
\providecommand{\e}{{\mathrm{e}}}
\providecommand{\vect}[1]{\overrightarrow{#1}}
\providecommand{\vectt}[1]{\overrightarrow{\mathrm{#1}}}
\providecommand{\vvt}[1]{\overrightarrow{\mathrm{#1}}}
\providecommand{\np}[2][]{{#2~#1}}
\providecommand{\vv}[1]{{\overrightarrow{#1}}}
\renewcommand{\d}{{\mathrm{d}}}
\providecommand{\degres}{{^\circ}}
\providecommand{\degre}{{^\circ}}
\providecommand{\blue}{\color{blue}}
\providecommand{\red}{\color{red}}
\providecommand{\enskip}{\quad}
\providecommand{\llbracket}{\lbrack\!\lbrack}
\providecommand{\rrbracket}{\rbrack\!\rbrack}
\providecommand{\euro}{{\text{€}}}
\providecommand{\no}{{\text{n°}}}
\providecommand{\interval}[2]{{[#1;#2]}}
\providecommand{\intervFO}[2]{{[#1;#2[}}
\providecommand{\dotfill}{\ldots}

\providecommand{\ds}{\displaystyle}
\providecommand{\textding}[1]{#1}
\providecommand{\Oij}{\(\mathrm{O};\vec{\imath},\vec{\jmath}\)}
\providecommand{\Oijk}{\(\mathrm{O};\vec{\imath},\vec{\jmath},\vec{k}\)}
\providecommand{\Ouv}{\(\mathrm{O};\vec{u},\vec{v}\)}
\providecommand{\intervFF}[2]{\ensuremath{[#1;#2]}}
\renewcommand{\euro}{\ifmmode\text{\texteuro}\else\texteuro\fi}
\DeclareUnicodeCharacter{00B0}{\ifmmode^{\circ}\else\textdegree\fi}
\DeclareUnicodeCharacter{20AC}{\ifmmode\text{\texteuro}\else\texteuro\fi}
% Énoncés : Éducation nationale. Transcription et corrigés : APMEP, Denis Vergès et contributeurs. Découpage et classement : CoopMaths, Sébastien Lozano, Éric Elter et contributeurs. Adaptation de la présentation : Maths Explorer.
% CC BY-SA : https://coopmaths.fr/www/about/
% https://coopmaths.fr/alea/static/bac/2021/tex/bac_2021_06_asie_9.tex
\begin{document}
\noindent\textit{APMEP / CoopMaths. CC BY-SA. Présentation adaptée par Maths Explorer.}\par\medskip
\medskip

\begin{tabular}{|l|}\hline
\textbf{Principaux domaines abordés}\\
-- Suites\\
-- Équations différentielles\\ \hline
\end{tabular}

\medskip

Dans cet exercice, on s'intéresse à la croissance du bambou Moso de taille maximale 20 mètres. 

Le modèle de croissance de Ludwig von Bertalanffy suppose que la vitesse de croissance pour un tel bambou est proportionnelle à l'écart entre sa taille et la taille maximale.

\bigskip

\textbf{Partie I : modèle discret}

\medskip

Dans cette partie, on observe un bambou de taille initiale $1$~mètre.

Pour tout entier naturel $n$, on note $u_n$ la taille, en mètre, du bambou $n$ jours après le début de l'observation. On a ainsi $u_0 = 1$.

Le modèle de von Bertalanffy pour la croissance du bambou entre deux jours consécutifs se traduit par l'égalité :

\[u_{n+1} = u_n + 0,05\left(20 - u_n\right)\,  \text{pour tout entier naturel}\,  n.\]

\medskip

\begin{enumerate}
\item Vérifier que $u_1 = 1,95$.
\item 
	\begin{enumerate}
		\item Montrer que pour tout entier naturel $n$,\,  $u_{n+1} = 0,95u_n + 1$.
		\item On pose pour tout entier naturel $n$,\,  $v_n = 20 - u_n$. 
		
Démontrer que la suite $\left(v_n\right)$ est une suite géométrique dont on précisera le terme initial $v_0$ et la raison.
		\item En déduire que, pour tout entier naturel $n$,\,  $u_n = 20 - 19 \times 0,95^n$.
	\end{enumerate}
\item Déterminer la limite de la suite $\left(u_n\right)$.
\end{enumerate}

\bigskip

\textbf{Partie II : modèle continu}

\medskip

Dans cette partie, on souhaite modéliser la taille du même bambou Moso par une fonction donnant sa taille, en mètre, en fonction du temps $t$ exprimé en jour. 

D'après le modèle de von Bertalanffy, cette fonction est solution de l'équation différentielle

\[(E) \qquad y' = 0,05(20 - y)\]

où $y$ désigne une fonction de la variable $t$, définie et dérivable sur $[0~;~+\infty[$ et $y'$ désigne sa fonction dérivée.

Soit la fonction $L$ définie sur l'intervalle $[0~;~+\infty[$ par 

\[L(t) = 20 - 19\e^{-0,05t}.\]

\smallskip

\begin{enumerate}
\item Vérifier que la fonction $L$ est une solution de $(E)$ et qu'on a également $L(0) = 1$.
\item On prend cette fonction $L$ comme modèle et on admet que, si on note $L'$ sa fonction dérivée,
$L'(t)$ représente la vitesse de croissance du bambou à l'instant $t$.
	\begin{enumerate}
		\item Comparer $L'(0)$ et $L'(5)$.
		\item Calculer la limite de la fonction dérivée $L'$ en $+\infty$. 
		
Ce résultat est-il en cohérence avec la description du modèle de croissance exposé au début de l'exercice ?
	\end{enumerate}
\end{enumerate}

\bigskip

\end{document}
